\section{Finite locally free morphisms}
\label{section-finite-locally-free}

\noindent
In many papers the authors use finite flat morphisms when they really mean
finite locally free morphisms. The reason is that if the base is locally
Noetherian then this is the same thing. But in general it is not, see
Exercises, Exercise \ref{exercises-exercise-flat-not-projective}.

\begin{definition}
\label{definition-finite-locally-free}
Let $f : X \to S$ be a morphism of schemes.
We say $f$ is {\it finite locally free} if $f$ is
affine and $f_*\mathcal{O}_X$ is a finite locally
free $\mathcal{O}_S$-module. In this case we say $f$ is
has {\it rank} or {\it degree} $d$
if the sheaf $f_*\mathcal{O}_X$ is finite locally free of degree $d$.
\end{definition}

\noindent
Note that if $f : X \to S$ is finite locally free then $S$ is the disjoint
union of open and closed subschemes $S_d$ such that $f^{-1}(S_d) \to S_d$
is finite locally free of degree $d$.

\begin{lemma}
\label{lemma-finite-flat}
Let $f : X \to S$ be a morphism of schemes.
The following are equivalent:
\begin{enumerate}
\item $f$ is finite locally free,
\item $f$ is finite, flat, and locally of finite presentation.
\end{enumerate}
If $S$ is locally Noetherian these are also equivalent to
\begin{enumerate}
\item[(3)] $f$ is finite and flat.
\end{enumerate}
\end{lemma}

\begin{proof}
Let $V \subset S$ be affine open. In all three cases
the morphism is affine hence $f^{-1}(V)$ is affine.
Thus we may write $V = \Spec(R)$ and $f^{-1}(V) = \Spec(A)$
for some $R$-algebra $A$.
Assume (1). This means we can cover $S$ by affine opens
$V = \Spec(R)$ such that $A$ is finite free as an $R$-module.
Then $R \to A$ is of finite presentation by
Algebra, Lemma \ref{algebra-lemma-finite-finite-type}.
Thus (2) holds. Conversely, assume (2).
For every affine open $V = \Spec(R)$ of $S$ the ring map
$R \to A$ is finite and of finite presentation and $A$ is flat
as an $R$-module. By
Algebra, Lemma \ref{algebra-lemma-finite-finitely-presented-extension}
we see that $A$ is finitely presented as an $R$-module.
Thus Algebra, Lemma \ref{algebra-lemma-finite-projective}
implies $A$ is finite locally free. Thus (1) holds.
The Noetherian case follows as a finite module over a Noetherian ring
is a finitely presented module, see Algebra,
Lemma \ref{algebra-lemma-Noetherian-finite-type-is-finite-presentation}.
\end{proof}

\begin{lemma}
\label{lemma-composition-finite-locally-free}
A composition of finite locally free morphisms is finite locally free.
\end{lemma}

\begin{proof}
Omitted.
\end{proof}

\begin{lemma}
\label{lemma-base-change-finite-locally-free}
A base change of a finite locally free morphism is finite locally free.
\end{lemma}

\begin{proof}
Omitted.
\end{proof}

\begin{lemma}
\label{lemma-finite-locally-free}
Let $f : X \to S$ be a finite locally free morphism of schemes.
There exists a disjoint union decomposition
$S = \coprod_{d \geq 0} S_d$ by open and closed subschemes
such that setting $X_d = f^{-1}(S_d)$ the restrictions
$f|_{X_d}$ are finite locally free morphisms $X_d \to S_d$
of degree $d$.
\end{lemma}

\begin{proof}
This is true because a finite locally free sheaf locally has
a well defined rank. Details omitted.
\end{proof}

\begin{lemma}
\label{lemma-massage-finite}
Let $f : Y \to X$ be a finite morphism with $X$ affine.
There exists a diagram
$$
\xymatrix{
Z' \ar[rd] &
Y' \ar[l]^i \ar[d] \ar[r] &
Y \ar[d] \\
 & X' \ar[r] & X
}
$$
where
\begin{enumerate}
\item $Y' \to Y$ and $X' \to X$ are surjective finite locally free,
\item $Y' = X' \times_X Y$,
\item $i : Y' \to Z'$ is a closed immersion,
\item $Z' \to X'$ is finite locally free, and
\item $Z' = \bigcup_{j = 1, \ldots, m} Z'_j$ is a (set theoretic)
finite union of closed subschemes, each of which maps isomorphically
to $X'$.
\end{enumerate}
\end{lemma}

\begin{proof}
Write $X = \Spec(A)$ and $Y = \Spec(B)$. See also
More on Algebra, Section \ref{more-algebra-section-descent-flatness-integral}.
Let $x_1, \ldots, x_n \in B$ be generators of $B$ over $A$.
For each $i$ we can choose a monic polynomial $P_i(T) \in A[T]$
such that $P(x_i) = 0$ in $B$. By
Algebra, Lemma \ref{algebra-lemma-adjoin-roots}
(applied $n$ times) there exists a finite locally free ring
extension $A \subset A'$ such that each $P_i$ splits completely:
$$
P_i(T) = \prod\nolimits_{k = 1, \ldots, d_i} (T - \alpha_{ik})
$$
for certain $\alpha_{ik} \in A'$. Set
$$
C = A'[T_1, \ldots, T_n]/(P_1(T_1), \ldots, P_n(T_n))
$$
and $B' = A' \otimes_A B$. The map $C \to B'$, $T_i \mapsto 1 \otimes x_i$
is an $A'$-algebra surjection. Setting $X' = \Spec(A')$,
$Y' = \Spec(B')$ and $Z' = \Spec(C)$ we see that
(1) -- (4) hold. Part (5) holds because set theoretically
$\Spec(C)$ is the union of the closed subschemes
cut out by the ideals
$$
(T_1 - \alpha_{1k_1}, T_2 - \alpha_{2k_2}, \ldots, T_n - \alpha_{nk_n})
$$
for any $1 \leq k_i \leq d_i$.
\end{proof}

\noindent
The following lemma is stated in the correct generality
in Lemma \ref{lemma-image-nowhere-dense-quasi-finite} below.

\begin{lemma}
\label{lemma-image-nowhere-dense-finite}
Let $f : Y \to X$ be a finite morphism of schemes.
Let $T \subset Y$ be a closed nowhere dense subset of $Y$.
Then $f(T) \subset X$ is a closed nowhere dense subset of $X$.
\end{lemma}

\begin{proof}
By Lemma \ref{lemma-finite-proper} we know that $f(T) \subset X$ is closed.
Let $X = \bigcup X_i$ be an affine covering.
Since $T$ is nowhere dense in $Y$, we see that also $T \cap f^{-1}(X_i)$
is nowhere dense in $f^{-1}(X_i)$. Hence if we can prove the theorem in the
affine case, then we see that $f(T) \cap X_i$ is nowhere dense.
This then implies that $T$ is nowhere dense in $X$ by
Topology, Lemma \ref{topology-lemma-nowhere-dense-local}.

\medskip\noindent
Assume $X$ is affine. Choose a diagram
$$
\xymatrix{
Z' \ar[rd] &
Y' \ar[l]^i \ar[d]^{f'} \ar[r]_a &
Y \ar[d]^f \\
 & X' \ar[r]^b & X
}
$$
as in Lemma \ref{lemma-massage-finite}. The morphisms $a$, $b$ are
open since they are finite locally free
(Lemmas \ref{lemma-finite-flat} and \ref{lemma-fppf-open}).
Hence $T' = a^{-1}(T)$ is nowhere dense, see
Topology, Lemma \ref{topology-lemma-open-inverse-image-closed-nowhere-dense}.
The morphism $b$ is surjective and open.
Hence, if we can prove $f'(T') = b^{-1}(f(T))$ is
nowhere dense, then $f(T)$ is nowhere dense, see
Topology, Lemma \ref{topology-lemma-open-inverse-image-closed-nowhere-dense}.
As $i$ is a closed immersion, by
Topology, Lemma \ref{topology-lemma-closed-image-nowhere-dense}
we see that $i(T') \subset Z'$ is closed and nowhere dense.
Thus we have reduced the problem to the case discussed
in the following paragraph.

\medskip\noindent
Assume that $Y = \bigcup_{i = 1, \ldots, n} Y_i$ is a finite union of
closed subsets, each mapping isomorphically to $X$. Consider
$T_i = Y_i \cap T$. If each of the $T_i$ is nowhere dense in $Y_i$,
then each $f(T_i)$ is nowhere dense in $X$ as $Y_i \to X$ is an isomorphism.
Hence $f(T) = f(T_i)$ is a finite union of nowhere dense closed
subsets of $X$ and we win, see
Topology, Lemma \ref{topology-lemma-nowhere-dense}.
Suppose not, say $T_1$ contains a nonempty open $V \subset Y_1$.
We are going to show this leads to a contradiction.
Consider $Y_2 \cap V \subset V$. This is either
a proper closed subset, or equal to $V$. In the first case we replace
$V$ by $V \setminus V \cap Y_2$, so $V \subset T_1$ is open in $Y_1$ and
does not meet $Y_2$. In the second case we have
$V \subset Y_1 \cap Y_2$ is open in both $Y_1$ and $Y_2$.
Repeat sequentially with $i = 3, \ldots, n$. The result is a disjoint
union decomposition
$$
\{1, \ldots, n\} = I_1 \amalg I_2, \quad 1 \in I_1
$$
and an open $V$ of $Y_1$ contained in $T_1$ such that $V \subset Y_i$
for $i \in I_1$ and $V \cap Y_i = \emptyset$ for $i \in I_2$. Set
$U = f(V)$. This is an open of $X$ since $f|_{Y_1} : Y_1 \to X$ is
an isomorphism. Then
$$
f^{-1}(U) = V\ \amalg\ \bigcup\nolimits_{i \in I_2} (Y_i \cap f^{-1}(U))
$$
As $\bigcup_{i \in I_2} Y_i$ is closed, this implies that
$V \subset f^{-1}(U)$ is open, hence $V \subset Y$ is open.
This contradicts the assumption that $T$ is nowhere dense in $Y$, as desired.
\end{proof}







